\section*{Chapter \ref{ch:ll:pre-trade-risk-and-kill-switches} --- Pre-Trade Risk and Kill Switches}

\begin{solution}{exo:ll:pre-trade-risk-and-kill-switches:1}
The band is $250\,000 \times 500 / 10\,000 = 12\,500$: buys up to 262\,500 and sells down to 237\,500 pass. The buy at 26.2500
and the sell at 23.7500 pass (they are on the boundary); the buy at 26.2600 and the sell at 23.7400 are refused.
\end{solution}

\begin{solution}{exo:ll:pre-trade-risk-and-kill-switches:2}
The worst-case long exposure is $2\,000 + 1\,500 + 1\,000 = 4\,500$, whether or not an order has been acknowledged; a buy of
600 would make 5\,100: refused. The largest buy accepted is 500.
\end{solution}

\begin{solution}{exo:ll:pre-trade-risk-and-kill-switches:3}
$4 \times 10^6 / 2\,700~\text{s} \approx 1\,481$ executions a second, and $397 \times 10^6 / 4 \times 10^6 \approx 99$ shares per
execution: small orders, many of them.
\end{solution}

\begin{solution}{exo:ll:pre-trade-risk-and-kill-switches:4}
One: the other fourteen repeat it within the window. A runaway that varied its quantity by a share, or its price by a tick,
from one order to the next would pass them all; the duplicate check catches loops that repeat themselves, not loops in general.
\end{solution}

\begin{solution}{exo:ll:pre-trade-risk-and-kill-switches:5}
Until two seconds after the last \omterm{def:ll:protocols-i-fix:session}{heartbeat}, which came at most \qty{400}{\milli\second} before $t$: between $t + 1.6$ and
$t + 2$ seconds. A gate that kept trusting its last limits until the service returned would fail open: the limits could have
been tightened, or the service could have stopped because something is wrong.
\end{solution}

\begin{solution}{exo:ll:pre-trade-risk-and-kill-switches:6}
About 3\% of the budget. A separate process on another core needs a round trip through rings, about \qty{200}{\nano\second}
at the median (chapter~12), eight times the check, plus a decision about what to do when it does not answer.
\end{solution}

\begin{solution}{exo:ll:pre-trade-risk-and-kill-switches:7}
The loss is updated on every fill (realised) and every reference price (unrealised), in the gate's report path, not in the
check; when it falls below the limit the gate kills the strategy (block and cancel). It is not a check on the order, since an
order does not change the loss; it is a monitor whose action is a kill. Test it with the runaway against a falling market.
\end{solution}

\begin{solution}{exo:ll:pre-trade-risk-and-kill-switches:8}
It is a post-trade monitor, not a pre-trade check: the orders it sees are already on their way, and the kill arrives after a
hop, while the runaway keeps sending (1\,500 orders a second leave in the time a person reads an alert). If the process stops,
nothing is checked and nothing says so: the design fails open. The check belongs on the path, before the gateway; a copy for
monitoring is a second line.
\end{solution}

\begin{solution}{pb:ll:pre-trade-risk-and-kill-switches:1}
\begin{enumerate}
\item $4 \times 10^6$ executions in $45 \times 60 = 2\,700$ seconds: about 1\,481 a second, so 1\,500 matches.
\item 750\,000 shares in five seconds; at 1\,500 orders of 100 shares a second, $1\,500 \times 100 \times 2\,700 = 405$ million
in 45 minutes, against the incident's 397 million.
\item 7\,500 fills make 749 ticks of 0.01 on a price of 100.00: about 7.5\%.
\item The collar is centred on the last trade, which the runaway moves; each child order is priced five ticks above it, well
inside a 5\% band. A collar catches prices far from the market, not a market pushed by the firm's own orders.
\item Each child order is 100 shares, about \$10\,000: far below any size or notional limit set for ordinary orders.
\item The bucket starts full with 50 tokens, which the first bursts spend in about \qty{40}{\milli\second}; afterwards the
runaway sends at the bucket's rate, 500 orders a second.
\item $500 \times 100 \times 2\,700 = 135$ million shares.
\item One order per window: once \qty{100}{\milli\second} have passed, the same order is no longer a duplicate, so about eleven
a second go through. A runaway that changed its quantity by one share each time would defeat the check.
\item 200 orders, 20\,000 shares, with the first refusal after \qty{130}{\milli\second}.
\item Counting only fills ignores the orders in flight: the last burst leaves before the fills of its first orders return (a
round trip of \qty{70}{\micro\second} against a burst of \qty{15}{\micro\second}), so ten more orders, 1\,000 shares, go out.
\item With more orders in flight: larger bursts, larger orders, a higher rate, or a slow report path (the incident's position
monitor lagged in high volume).
\item \$2.5 million at about \$100 a share is about 25\,000 shares; the notional is counted at the order's price, five ticks
above a rising last trade, so the threshold is reached a little earlier, at 24\,900.
\item \textbf{Named result.} At the incident's rate of about 1\,500 orders a second: with no check, or a collar, or size and
notional limits per order, 7\,500 orders and 750\,000 shares in five seconds, 405 million shares carried over 45 minutes (the
incident: 397 million); with a throttle at 500 a second, 135 million; with a duplicate check, 3.0 million; with a position limit
of 20\,000 counted in flight, 200 orders and 20\,000 shares after \qty{130}{\milli\second} (21\,000 counting fills only); with a
capital threshold of \$2.5 million, 249 orders and 24\,900 shares.
\item It blocks the strategy instead of refusing its orders one by one (7\,300 refusals in five seconds without it), cancels
its open orders, and stops it in every instrument at once.
\item A per-instrument limit stops the runaway in each instrument separately: across 154 instruments, 154 times the limit. A
firm-wide capital threshold, or a kill switch at the first breach, is immune to that.
\item The collar row is the 9.5\% price cap; the position rows are the position limits that did not count outstanding
orders (fills only) and the fix (in flight); the capital row is the missing firm-wide threshold; the kill switch is the missing
procedure to halt the router; the duplicate and throttle rows are controls on the router's own output.
\item Because the moment a gate cannot check is the moment something is wrong; passing orders then gives no protection when
it is needed.
\item It sees executions after they happen, through a path that can lag, and a person must act on it: at 1\,500 orders a
second, every second of delay is 150\,000 shares.
\item Strategies that legitimately repeat orders: a quote refreshed at the same price and size, an order re-sent after a
cancel, slices of an execution algorithm, retries of immediate-or-cancel orders.
\item A limit on what the firm holds or could hold, counted in the worst case and linked to an automatic kill.
\end{enumerate}
\end{solution}

\begin{solution}{iq:ll:pre-trade-risk-and-kill-switches:1}
Price collar, maximum size and notional per order, position limits counting orders in flight, capital or gross notional
thresholds at desk and firm level, order and message rates, duplicate detection, open-order counts, stale-data checks, and a kill
switch. The ones that catch a runaway are those on accumulated exposure (position, capital) and on rates; per-order checks
catch fat fingers.
\iqlookfor{the list, and the distinction between per-order and accumulated checks.}
\end{solution}

\begin{solution}{iq:ll:pre-trade-risk-and-kill-switches:2}
Precompute everything (collar bands, limits per instrument in dense arrays indexed by a small integer), keep exposures as
running sums updated by reports, compute every condition into a bit mask without branches, and refuse on a nonzero mask. Measure
it on both outcomes.
\iqlookfor{precomputation, running sums, branch-free evaluation, measurement.}
\end{solution}

\begin{solution}{iq:ll:pre-trade-risk-and-kill-switches:3}
It keeps trading on the last limits only for their stated maximum age, then refuses every order (fail closed) and alerts. The
limits may have been tightened, and the silence may be a symptom of a larger failure.
\iqlookfor{fail closed, with a bounded grace period.}
\end{solution}

\begin{solution}{iq:ll:pre-trade-risk-and-kill-switches:4}
Block new orders, cancel open ones (the cancels are in flight until confirmed), optionally propose flattening orders for a
person to confirm. The gate itself on a hard breach, the risk desk and the trader can pull it; only an authorised person, after
understanding the cause, undoes it.
\iqlookfor{block, cancel, flatten; automatic pull, manual release.}
\end{solution}

\begin{solution}{iq:ll:pre-trade-risk-and-kill-switches:5}
Publish a complete, versioned snapshot; the gate installs it between two orders, never during a check, and records the version
with every decision. Checks against the new limits apply to the exposure already accumulated.
\iqlookfor{atomic snapshots between events, versioned.}
\end{solution}

\begin{solution}{iq:ll:pre-trade-risk-and-kill-switches:6}
Controls on the output of every order-generating component, before the venue: worst-case position and capital limits linked to
an automatic kill, rate and duplicate checks, collars; fail closed; post-trade monitoring as a second line. Test by replaying a
synthetic runaway through the whole path, one control at a time, and asserting the orders and position at which each trips.
\iqlookfor{requirements traced to the incident, and a runaway replay as the test.}
\end{solution}
